\documentclass[11pt]{article}

\usepackage[margin=1in]{geometry}
\usepackage{amsmath,amssymb,amsthm,mathtools}
\usepackage{bm}
\usepackage{enumitem}
\usepackage{hyperref}

\title{A Primer on the Implemented Korea Spatial GE Model}
\author{}
\date{\today}

\newcommand{\R}{\mathcal{N}}
\newcommand{\Sct}{\mathcal{J}}
\newcommand{\E}{\mathbb{E}}

\begin{document}
\maketitle

\section{Purpose}
This note documents the quantitative model implemented in the \texttt{korea\_growth}
package. It is a background primer rather than a full theory paper: the goal is to
state the objects, equilibrium conditions, and numerical solution method that are
actually used by the code.

The implementation solves a sequence of static spatial equilibria. The only dynamic
state is lagged population, so period $t$ equilibrium takes $L_{t-1}$ as given and
returns $(w_t,r_t,P_t,E_t,L_t,\bar\tau_t,\bar\pi_t)$.

\section{Environment}
Time is discrete, $t = 1,\dots,T$. Regions are indexed by $o,d \in \R$ and sectors by
$j,k \in \Sct$. Let $N = |\R|$ and $J = |\Sct|$.

The exogenous inputs in period $t$ are:
\begin{itemize}[leftmargin=1.5em]
\item regional amenities $\bar V_{d,t}$,
\item domestic trade costs $\tau_{odj,t} \ge 1$ and import costs $\tilde\tau_{oj,t} \ge 1$,
\item migration costs $\delta_{od,t} \ge 1$,
\item sectoral productivities $A_{oj,t}$,
\item land shares $\beta_{oj,t}$ and value-added shares $\gamma_{oj,t}$,
\item input-output shares $\gamma^{IO}_{ojk,t}$,
\item fixed costs $F_{oj,t}$, export fixed costs $\tilde F_{oj,t}$, and adoption fixed costs
      $\breve F_{oj,t}$,
\item subsidy rates $s_{oj,t}$,
\item foreign demand shifters $\tilde D_{j,t}$ and foreign prices $\tilde p_{j,t}$,
\item land or structures supply $H_{o,t}$,
\item the mass of potential firms $M_{j,t}$.
\end{itemize}

For the agriculture sector there is an alternative mechanized technology with
$(\breve\beta_{oj,t}, \breve\gamma_{oj,t}, \breve\gamma^{IO}_{ojk,t})$ and efficiency gain
$\xi > 1$.

The structural parameters are:
\[
\sigma > 1,\quad \theta > 0,\quad \kappa > 1,\quad \rho_j \in (0,1),\quad
\iota \in (-1,0),\quad \nu > 0,\quad \eta \in \mathbb{R},
\]
together with expenditure parameters $\alpha_j$ and $v_j$ satisfying
$\sum_j \alpha_j = 1$ and $\sum_j v_j = 0$.

\section{Households}
\subsection{Disposable income and indirect utility}
Per-capita disposable income in region $d$ is
\begin{equation}
  y^{pc}_{d,t} = (1 - \bar\tau_t + \bar\pi_t) w_{d,t},
\end{equation}
where $\bar\tau_t$ is the proportional tax needed to finance subsidies and infrastructure
and $\bar\pi_t$ is aggregate profit income rebated in wage units.

Given sectoral price indices $P_{djt}$, define the composite price index
\begin{equation}
  \mathcal{P}_{d,t} = \prod_{j \in \Sct} P_{djt}^{\alpha_j}.
\end{equation}
The implemented indirect utility is
\begin{equation}
  V_{d,t} =
  \begin{cases}
    \dfrac{1}{\eta} \left( \dfrac{y^{pc}_{d,t}}{\mathcal{P}_{d,t}} \right)^{\eta}
    - \sum_j v_j \log P_{djt}, & \eta \ne 0, \\[1em]
    \log\!\left(\dfrac{y^{pc}_{d,t}}{\mathcal{P}_{d,t}}\right)
    - \sum_j v_j \log P_{djt}, & \eta = 0.
  \end{cases}
\end{equation}

\subsection{Migration}
Lagged population $L_{d,t-1}$ enters destination utility through congestion:
\begin{equation}
  g_{d,t} = L_{d,t-1}^{\iota}.
\end{equation}
The common component of utility from moving to $d$ is
\begin{equation}
  U^{com}_{d,t} = \bar V_{d,t} g_{d,t} V_{d,t}.
\end{equation}
An agent starting in origin $o$ values destination $d$ at
\begin{equation}
  U_{od,t} = \frac{U^{com}_{d,t}}{\delta_{od,t}}.
\end{equation}
Migration shares are multinomial-logit in levels:
\begin{equation}
  \mu_{od,t} = \frac{U_{od,t}^{\nu}}{\sum_{d' \in \R} U_{od',t}^{\nu}}.
\end{equation}
Population evolves according to
\begin{equation}
  L_{d,t} = \sum_{o \in \R} \mu_{od,t} L_{o,t-1}.
\end{equation}

\subsection{Non-homothetic final demand}
Expenditure shares in destination $d$ are
\begin{equation}
  \psi_{dj,t}
  = \alpha_j + v_j \left( \frac{y^{pc}_{d,t}}{\mathcal{P}_{d,t}} \right)^{-\eta}.
\end{equation}
Final consumption expenditure is
\begin{equation}
  E^{F}_{dj,t} = \psi_{dj,t} y^{pc}_{d,t} L_{d,t}.
\end{equation}

\section{Production and firm heterogeneity}
\subsection{Unit costs and agglomeration}
Agglomeration is sector specific:
\begin{equation}
  f_{oj,t} = L_{o,t-1}^{\rho_j}.
\end{equation}

For the baseline technology, unit cost in region $o$ and sector $j$ is
\begin{equation}
  c_{oj,t}
  =
  \left[
    \frac{1}{\gamma_{oj,t}}
    \left( \frac{r_{o,t}}{\beta_{oj,t}} \right)^{\beta_{oj,t}}
    \left( \frac{w_{o,t}}{1-\beta_{oj,t}} \right)^{1-\beta_{oj,t}}
  \right]^{\gamma_{oj,t}}
  \prod_{k \in \Sct}
  \left( \frac{P_{ok,t}}{\gamma^{IO}_{ojk,t}} \right)^{\gamma^{IO}_{ojk,t}}.
\end{equation}
The mechanized agricultural technology has an analogous unit cost
$\breve c_{oj,t}$ built from $(\breve\beta,\breve\gamma,\breve\gamma^{IO})$.

\subsection{Productivity draws}
Potential entrants draw idiosyncratic productivity $\phi \in [1,\kappa]$ from a bounded
Pareto distribution with cdf
\begin{equation}
  G(\phi) = \frac{1 - \phi^{-\theta}}{1 - \kappa^{-\theta}},
  \qquad \phi \in [1,\kappa],
\end{equation}
and density
\begin{equation}
  g(\phi) = \frac{\theta \phi^{-\theta-1}}{1 - \kappa^{-\theta}}.
\end{equation}
The main integral used by the code is
\begin{equation}
  I(a,b) = \int_a^b \phi^{\sigma-1} g(\phi)\, d\phi.
\end{equation}

\subsection{Demand shifter}
Given destination expenditures and prices, origin $o$ faces sector-$j$ demand shifter
\begin{equation}
  S_{oj,t}
  = \sum_{d \in \R} \tau_{odj,t}^{1-\sigma} P_{djt}^{\sigma-1} E_{dj,t}.
\end{equation}

\subsection{Entry, adoption, and export cutoffs}
The zero-profit entry cutoff is implemented as
\begin{equation}
  \bar\phi_{oj,t}
  =
  \max\left\{
    1,\,
    \frac{\sigma}{\sigma-1}
    \frac{1-s_{oj,t}}{A_{oj,t} f_{oj,t}}
    \sigma^{\frac{1}{\sigma-1}}
    c_{oj,t}
    \left( \frac{F_{oj,t}}{S_{oj,t}} \right)^{\frac{1}{\sigma-1}}
  \right\}.
\end{equation}

For agriculture, the adoption cutoff uses the mechanization benefit term
\begin{equation}
  B_{oj,t}
  =
  \left[
    \left( \xi \frac{c_{oj,t}}{\breve c_{oj,t}} \right)^{\sigma-1} - 1
  \right]^{\frac{1}{1-\sigma}},
\end{equation}
whenever the expression in brackets is positive. Then
\begin{equation}
  \breve\phi_{oj,t}
  =
  \max\left\{
    1,\,
    \frac{\sigma}{\sigma-1}
    \frac{1-s_{oj,t}}{A_{oj,t} f_{oj,t}}
    \sigma^{\frac{1}{\sigma-1}}
    c_{oj,t} B_{oj,t}
    \left( \frac{\breve F_{oj,t}}{S_{oj,t}} \right)^{\frac{1}{\sigma-1}}
  \right\}.
\end{equation}
If the benefit term is non-positive, the code sets $\breve\phi_{oj,t} = \infty$ and
mechanized adoption never occurs.

The export cutoff is
\begin{equation}
  \tilde\phi_{oj,t}
  =
  \max\left\{
    1,\,
    \frac{\sigma}{\sigma-1}
    \frac{1-s_{oj,t}}{A_{oj,t} f_{oj,t}}
    \sigma^{\frac{1}{\sigma-1}}
    c_{oj,t}
    \left(
      \frac{\tilde F_{oj,t}}
      {\tilde\tau_{oj,t}^{1-\sigma} \tilde D_{j,t}}
    \right)^{\frac{1}{\sigma-1}}
    \times \chi_{oj,t}
  \right\},
\end{equation}
where $\chi_{oj,t} = (\breve c_{oj,t}/c_{oj,t}) / \xi$ for agriculture and
$\chi_{oj,t} = 1$ otherwise.

\subsection{Productivity aggregates}
Define
\begin{align}
  Z_{oj,t} &= A_{oj,t} f_{oj,t} I(\bar\phi_{oj,t}, \bar b_{oj,t})^{\frac{1}{\sigma-1}}, \\
  \breve Z_{oj,t} &= A_{oj,t} f_{oj,t} I(\breve\phi_{oj,t}, \kappa)^{\frac{1}{\sigma-1}}, \\
  \tilde Z_{oj,t} &= A_{oj,t} f_{oj,t} \chi^x_j I(\tilde\phi_{oj,t}, \kappa)^{\frac{1}{\sigma-1}},
\end{align}
where $\bar b_{oj,t} = \min\{\breve\phi_{oj,t}, \kappa\}$ in agriculture and
$\bar b_{oj,t} = \kappa$ otherwise, and $\chi^x_j = \xi$ for agriculture and $1$ otherwise.

\section{Sectoral price indices and revenues}
\subsection{Price indices}
Let
\begin{equation}
  \Omega_{oj,t}
  =
  M_{j,t}
  \left( \frac{\sigma}{\sigma-1}(1-s_{oj,t}) \right)^{1-\sigma}.
\end{equation}
The baseline and mechanized domestic supply terms are
\begin{align}
  B^{trad}_{oj,t} &= \Omega_{oj,t} \left( \frac{Z_{oj,t}}{c_{oj,t}} \right)^{\sigma-1}, \\
  B^{new}_{oj,t} &=
  \begin{cases}
    \Omega_{oj,t} \left( \dfrac{\breve Z_{oj,t}}{\breve c_{oj,t}} \right)^{\sigma-1},
    & \text{agriculture}, \\
    0, & \text{otherwise}.
  \end{cases}
\end{align}
Then the destination-$d$ price index in sector $j$ satisfies
\begin{equation}
  P_{djt}^{1-\sigma}
  =
  \sum_{o \in \R} \left( B^{trad}_{oj,t} + B^{new}_{oj,t} \right) \tau_{odj,t}^{1-\sigma}
  + (\tilde\tau_{dj,t} \tilde p_{j,t})^{1-\sigma}.
\end{equation}

\subsection{Domestic and foreign revenues}
Domestic revenues are
\begin{align}
  R^{trad}_{oj,t} &= B^{trad}_{oj,t} S_{oj,t}, \\
  R^{new}_{oj,t} &= B^{new}_{oj,t} S_{oj,t}, \\
  R^{dom}_{oj,t} &= R^{trad}_{oj,t} + R^{new}_{oj,t}.
\end{align}

Export revenue is
\begin{equation}
  R^{x}_{oj,t}
  =
  M_{j,t}
  \left(
    \frac{\sigma}{\sigma-1}
    (1-s_{oj,t}) \tilde\tau_{oj,t}
  \right)^{1-\sigma}
  \left( \frac{\tilde Z_{oj,t}}{c^x_{oj,t}} \right)^{\sigma-1}
  \tilde D_{j,t},
\end{equation}
where $c^x_{oj,t} = \breve c_{oj,t}$ for agriculture and $c^x_{oj,t} = c_{oj,t}$ otherwise.
Gross output is
\begin{equation}
  Y_{oj,t} = R^{dom}_{oj,t} + R^{x}_{oj,t}.
\end{equation}

The mass of active firms, exporters, and adopters is
\begin{align}
  N^{firm}_{oj,t} &= M_{j,t} [1 - G(\bar\phi_{oj,t})], \\
  N^{x}_{oj,t} &= M_{j,t} [1 - G(\tilde\phi_{oj,t})], \\
  N^{adopt}_{oj,t} &=
  \begin{cases}
    M_{j,t} [1 - G(\breve\phi_{oj,t})], & \text{agriculture}, \\
    0, & \text{otherwise}.
  \end{cases}
\end{align}

\section{Goods markets and factor markets}
\subsection{Expenditure system}
Under CES demand with constant markups, variable costs are a fixed share of revenues:
\begin{align}
  VC_{oj,t} &= \frac{\sigma-1}{\sigma}
  \left( R^{trad}_{oj,t} + R^{new}_{oj,t} + R^{x}_{oj,t} \right), \\
  VC^{trad}_{oj,t} &= \frac{\sigma-1}{\sigma} R^{trad}_{oj,t}, \\
  VC^{new}_{oj,t} &= \frac{\sigma-1}{\sigma}
  \left( R^{new}_{oj,t} + R^{x}_{oj,t} \right).
\end{align}

Local intermediate demand for sector $j$ in region $o$ is constructed from sectoral
variable costs:
\begin{equation}
  E^{M}_{oj,t} = \sum_{k \in \Sct} VC_{ok,t} \gamma^{IO}_{okj,t},
\end{equation}
except in agriculture, where the code replaces the baseline input share on the
agricultural output block with
\begin{equation}
  VC^{trad}_{oa,t} \gamma^{IO}_{oak,t}
  + VC^{new}_{oa,t} \breve\gamma^{IO}_{oak,t}.
\end{equation}
Total expenditure is
\begin{equation}
  E_{oj,t} = E^{F}_{oj,t} + E^{M}_{oj,t}.
\end{equation}

\subsection{Factor prices}
Labor income equals the labor share of variable costs and land income equals the land
share of variable costs. For non-agriculture,
\begin{align}
  w_{o,t} L_{o,t} &= \sum_j (1-\beta_{oj,t}) \gamma_{oj,t} VC_{oj,t}, \\
  r_{o,t} H_{o,t} &= \sum_j \beta_{oj,t} \gamma_{oj,t} VC_{oj,t}.
\end{align}
For agriculture, the code splits baseline and mechanized production:
\begin{align}
  w_{o,t} L_{o,t}
  &=
  (1-\beta_{oa,t}) \gamma_{oa,t} VC^{trad}_{oa,t}
  + (1-\breve\beta_{oa,t}) \breve\gamma_{oa,t} VC^{new}_{oa,t}
  + \sum_{j \ne a} (1-\beta_{oj,t}) \gamma_{oj,t} VC_{oj,t}, \\
  r_{o,t} H_{o,t}
  &=
  \beta_{oa,t} \gamma_{oa,t} VC^{trad}_{oa,t}
  + \breve\beta_{oa,t} \breve\gamma_{oa,t} VC^{new}_{oa,t}
  + \sum_{j \ne a} \beta_{oj,t} \gamma_{oj,t} VC_{oj,t}.
\end{align}

\section{Government budget and profit rebate}
Aggregate wage income is
\begin{equation}
  W_t = \sum_{o \in \R} w_{o,t} L_{o,t}.
\end{equation}
Aggregate profits are
\begin{equation}
  \Pi_t
  =
  \sum_{o \in \R} \sum_{j \in \Sct}
  \left[
    \frac{1}{\sigma} Y_{oj,t}
    - N^{firm}_{oj,t} F_{oj,t}
    - N^{x}_{oj,t} \tilde F_{oj,t}
    - N^{adopt}_{oj,t} \breve F_{oj,t}
  \right].
\end{equation}
The profit rebate is
\begin{equation}
  \bar\pi_t = \frac{\Pi_t}{W_t}.
\end{equation}

Subsidy and infrastructure spending imply
\begin{equation}
  \bar\tau_t
  =
  \frac{
    \sum_{o,j} \frac{\sigma-1}{\sigma} s_{oj,t} Y_{oj,t}
    + G^{infra}_t
  }{W_t},
\end{equation}
where $G^{infra}_t$ is exogenous road and building expenditure.

\section{Static equilibrium}
Given $L_{t-1}$ and exogenous period-$t$ fundamentals, a static equilibrium is a tuple
\[
  \left\{
    w_{o,t}, r_{o,t}, P_{oj,t}, E_{oj,t}, L_{o,t}, \bar\tau_t, \bar\pi_t
  \right\}_{o \in \R,\, j \in \Sct}
\]
such that:
\begin{enumerate}[leftmargin=1.5em]
\item household migration and final demand satisfy the equations in Section 2,
\item firm cutoffs, price indices, and revenues satisfy the equations in Sections 3 and 4,
\item expenditure balances satisfy $E_{oj,t} = E^{F}_{oj,t} + E^{M}_{oj,t}$,
\item factor markets satisfy the wage and land-rent equations in Section 5,
\item $\bar\pi_t$ and $\bar\tau_t$ satisfy the aggregate budget formulas in Section 6.
\end{enumerate}

The code treats the static equilibrium as a fixed point in
\[
  x_t = (w_t, r_t, P_t, E_t, \bar\tau_t, \bar\pi_t),
\]
with $L_t$ implied from $x_t$ and $L_{t-1}$.

\section{Dynamic path}
The full dynamic equilibrium is computed sequentially because the only endogenous
state variable is lagged population:
\begin{equation}
  L_{t-1} \longrightarrow \text{solve static equilibrium at } t \longrightarrow L_t.
\end{equation}
For $t=1$, the state is the initial population vector $L_0$. For $t \ge 2$, the
state is the solved population vector from period $t-1$.

\section{Numerical solution}
The solver uses a damped fixed-point iteration:
\begin{enumerate}[leftmargin=1.5em]
\item start from a positive guess for $(w,r,P,E,\bar\tau,\bar\pi)$,
\item compute the implied mapping $F(x_t)$ from the equilibrium equations,
\item update each component with geometric damping in logs,
\item stop when the maximum absolute log residual falls below tolerance.
\end{enumerate}

For example, wages are updated by
\begin{equation}
  w^{new} = \exp\left[(1-\lambda)\log w^{old} + \lambda \log w^{imp}\right],
\end{equation}
and analogously for $r$, $P$, $E$, $\bar\tau$, and $\bar\pi$.
Later periods are warm-started from the previous period's solution.

\section{Implementation restrictions}
The code enforces several restrictions that are useful to keep in mind when extending
the model:
\begin{itemize}[leftmargin=1.5em]
\item $L_0$ is normalized to sum to one.
\item $\gamma_{oj,t} + \sum_k \gamma^{IO}_{ojk,t} = 1$ for every $(o,j,t)$.
\item $\breve\gamma_{oj,t} + \sum_k \breve\gamma^{IO}_{ojk,t} = 1$ for every $(o,j,t)$.
\item $\breve\gamma_{oj,t}\breve\beta_{oj,t} = \gamma_{oj,t}\beta_{oj,t}$, so the
      land share inside value added is preserved across agricultural technologies.
\item Outside the heavy-manufacturing input, the mechanized agricultural input shares
      are restricted to equal the baseline input shares.
\item Subsidy rates satisfy $0 \le s_{oj,t} < 1$.
\end{itemize}

\section{Interpretation}
The model combines five ingredients:
\begin{enumerate}[leftmargin=1.5em]
\item non-homothetic demand and endogenous migration,
\item Melitz-style firm heterogeneity with bounded Pareto draws,
\item input-output linkages across sectors,
\item sector-specific agglomeration from lagged population,
\item policy wedges acting through productivity, fixed costs, subsidies, trade costs,
      migration costs, and foreign demand.
\end{enumerate}

That is the economic system solved by the repository. Any empirical exercise or policy
counterfactual run through the package is a perturbation of this equilibrium system.

\end{document}
